\begin{table}[H]
\centering
\small
\caption{GII y componentes en Argentina, India y Estados Unidos.}
\label{tab:gii_comparacion}
\renewcommand{\arraystretch}{1.25}
\begin{tabularx}{\textwidth}{>{\raggedright\arraybackslash}Xrrr}
\toprule
\textbf{Indicador} & \textbf{Argentina} & \textbf{India} & \textbf{Estados Unidos} \\
\midrule
Natalidad adolescente (nacimientos por 1.000 mujeres de 15--19 años) & 26,41 & 14,06 & 13,15 \\
Índice de desigualdad de género (GII, de 0 a 1) & 0,26 & 0,40 & 0,17 \\
Posición en el ranking del GII & 70 & 102 & 45 \\
Participación laboral de mujeres de 15 años o más (\%) & 53,16 & 35,08 & 57,33 \\
Participación laboral de varones de 15 años o más (\%) & 72,22 & 76,41 & 68,12 \\
Razón de mortalidad materna (muertes por 100.000 nacidos vivos) & 44,90 & 102,65 & 21,09 \\
Bancas parlamentarias ocupadas por mujeres (\%) & 43,77 & 14,82 & 28,22 \\
Bancas parlamentarias ocupadas por varones (\%) & 56,23 & 85,18 & 71,78 \\
Mujeres de 25 años o más con al menos parte de educación secundaria (\%) & 74,30 & 43,50 & 97,87 \\
Varones de 25 años o más con al menos parte de educación secundaria (\%) & 72,22 & 61,09 & 97,83 \\
\bottomrule
\end{tabularx}
\end{table}
